\documentclass[11pt]{article}

\usepackage[review]{acl}
\usepackage{times}
\usepackage{latexsym}
\usepackage[T1]{fontenc}
\usepackage[utf8]{inputenc}
\usepackage{microtype}
\usepackage{inconsolata}
\usepackage{graphicx}
\usepackage{amsmath}
\usepackage{array}

\title{Code-JEPA: Learning Code Representations in Latent Space}

\author{First Author \\
  Affiliation / Address line 1 \\
  \texttt{email@domain} \\\And
  Second Author \\
  Affiliation / Address line 1 \\
  \texttt{email@domain} \\}

\begin{document}
\maketitle

\begin{abstract}
Code-JEPA trains compact code encoders with joint embedding prediction over
semantics-preserving code views and hard negatives. Predictive and ranking losses
act in a projected latent space, leaving the pooled encoder embedding reusable
for downstream tasks. We sketch the method and a same-scale comparison protocol
against compact CodeBERT-, GraphCodeBERT-, TreeBERT-, and UniXcoder-style
baselines.
\end{abstract}

\section{Introduction}

Pretrained code encoders are usually optimized with token reconstruction,
natural-language alignment, or structure-aware auxiliary tasks. These objectives
learn useful representations, but they do not directly ask a model to predict
which semantics-preserving code variants should be close and which hard
negatives should be separated. Code-JEPA explores this question with a compact
joint embedding predictive architecture for code.

The main design issue is where to apply predictive and ranking pressure. If the
rank loss acts directly on the pooled encoder embedding, downstream tasks inherit
distortions introduced by hard-negative training. We instead keep the pooled
encoder representation \(h\) as the reusable embedding and train the predictive
objective in a projected latent space \(z\). This paper asks whether that
separation improves code representation learning under a small-model budget.

Our contributions are: (i) a Code-JEPA objective for transformed code views and
hard negatives, (ii) a projection-head design that protects downstream
embeddings from direct rank-loss distortion, and (iii) a same-scale experimental
protocol comparing Code-JEPA to compact CodeBERT-, GraphCodeBERT-, TreeBERT-,
and UniXcoder-style baselines.

\begin{figure*}[t!]
  \centering
  \setlength{\fboxsep}{4pt}
  \footnotesize
  \begin{tabular}{@{}c@{\hspace{0.5em}}c@{\hspace{0.5em}}c@{\hspace{0.5em}}c@{\hspace{0.5em}}c@{}}
    \fbox{\begin{minipage}{0.16\linewidth}\centering
      Anchor code\\[-1pt]\texttt{x\_a}
    \end{minipage}}
    &
    \fbox{\begin{minipage}{0.16\linewidth}\centering
      Positive view\\[-1pt]\texttt{x\_p}
    \end{minipage}}
    &
    \fbox{\begin{minipage}{0.16\linewidth}\centering
      Hard negative\\[-1pt]\texttt{x\_n}
    \end{minipage}}
    &
    \fbox{\begin{minipage}{0.18\linewidth}\centering
      Shared encoder\\[-1pt]\(f_\theta\)
    \end{minipage}}
    &
    \fbox{\begin{minipage}{0.18\linewidth}\centering
      Downstream tasks\\[-1pt]use \(h\)
    \end{minipage}}
    \\
    \multicolumn{3}{c}{\(\underbrace{\hspace{0.54\linewidth}}_{\text{semantics-preserving transforms and hard-negative mining}}\)}
    & \(\longrightarrow\)
    & \(\longrightarrow\)
  \end{tabular}

  \vspace{0.5em}
  \begin{tabular}{@{}c@{\hspace{0.45em}}c@{\hspace{0.45em}}c@{\hspace{0.45em}}c@{\hspace{0.45em}}c@{\hspace{0.45em}}c@{\hspace{0.45em}}c@{}}
    \fbox{\begin{minipage}{0.20\linewidth}\centering
      \(h_i=\mathrm{Pool}(f_\theta(x_i))\)\\
      reusable embedding
    \end{minipage}}
    \(\longrightarrow\)
    \fbox{\begin{minipage}{0.22\linewidth}\centering
      Projection head \(g_\phi\)\\
      Dense \(8H\), SwiGLU, RMSNorm
    \end{minipage}}
    \(\longrightarrow\)
    \fbox{\begin{minipage}{0.20\linewidth}\centering
      \(z_i=g_\phi(h_i)\)\\
      JEPA latent space
    \end{minipage}}
    \(\longrightarrow\)
    \fbox{\begin{minipage}{0.20\linewidth}\centering
      Predictor \(q_\psi\)\\
      \(\hat{z}_a=q_\psi(z_a)\)
    \end{minipage}}
  \end{tabular}

  \vspace{0.5em}
  \begin{equation}
    \begin{aligned}
      \mathcal{L}_{\mathrm{pred}} &= \|\hat{z}_a-\mathrm{sg}(z_p)\|_2^2,\\
      \mathcal{L}_{\mathrm{rank}} &= \max(0, m+\cos(\hat{z}_a,\mathrm{sg}(z_n))
          -\cos(\hat{z}_a,\mathrm{sg}(z_p))),\\
      \mathcal{L} &= \mathcal{L}_{\mathrm{pred}}
        +\lambda_r\mathcal{L}_{\mathrm{rank}}
        +\lambda_b\mathcal{L}_{\mathrm{inbatch}}
        +\lambda_s\mathcal{L}_{\mathrm{sigreg}} .
    \end{aligned}
  \end{equation}
  \caption{Code-JEPA architecture sketch. The reusable downstream embedding
  \(h\) is kept separate from the projected latent \(z\), where JEPA, rank,
  in-batch, and SIGReg losses are applied.}
  \label{fig:architecture}
\end{figure*}

\section{Method}

\paragraph{Code views.}
Given a Python function, we construct transformed views using formatting,
normalization, docstring/comment changes, and syntax-preserving edits. A training
triple consists of an anchor \(x_a\), a positive view \(x_p\), and a hard
negative \(x_n\) selected to be locally plausible but semantically farther.

\paragraph{Encoder and projection head.}
All models use a small UniXcoder-style RoBERTa backbone in the main experiments.
The encoder produces token states that are mean-pooled into \(h\). Code-JEPA then
uses a projection head \(g_\phi\) before applying predictive or ranking losses:
Dense \(8H\), split into two \(4H\) branches, SwiGLU, RMSNorm, and Dense \(H\).
This keeps \(h\) aligned with downstream use while allowing \(z\) to absorb
hard-negative pressure.

\paragraph{Objective.}
The predictor \(q_\psi\) maps the projected anchor \(z_a\) to \(\hat{z}_a\).
The positive and negative targets are stop-gradient latents. We combine a latent
prediction loss, a margin rank loss, an in-batch contrastive loss, and SIGReg
regularization. The rank weight is the main Code-JEPA knob; the matched control
sets it to zero.

\section{Experiments}

\paragraph{Pretraining data.}
Code-JEPA uses transformed CodeSearchNet Python views and triples. The raw
UniXcoder-style baseline uses raw CodeSearchNet Python function/docstring pairs
with masked language modeling and code-doc contrastive learning. All small-model
runs use the same tokenizer and comparable backbone scale.

\paragraph{Downstream tasks.}
We evaluate on POJ-104 clone retrieval with MAP@R and BigCloneBench clone
detection with F1 and ROC-AUC. Fine-tuning reads the pooled encoder embedding
\(h\), not the projection head, so the downstream protocol directly tests the
quality of the reusable representation.

\begin{table*}[t]
  \centering
  \small
  \begin{tabular}{lcccccc}
    \hline
    \textbf{Model} & \textbf{Scale} & \textbf{Pretraining data} & \textbf{Objective} &
    \textbf{Hard neg.} & \textbf{POJ MAP@R} & \textbf{BCB F1 / AUC} \\
    \hline
    Small CodeBERT-style & same & CodeSearchNet raw & MLM & no & TBD & TBD \\
    Small GraphCodeBERT-style & same & Code + graph signal & MLM + graph & partial & TBD & TBD \\
    Small TreeBERT-style & same & Code + AST signal & tree-aware MLM & no & TBD & TBD \\
    Small UniXcoder-style & same & CodeSearchNet raw & MLM + code-doc CL & no & 0.440 & 0.917 / 0.977 \\
    Code-JEPA & same & CodeSearchNet transforms & JEPA + rank + SIGReg & yes & TBD & TBD \\
    \hline
  \end{tabular}
  \caption{Planned same-scale comparison table. Values marked TBD should be
  filled only after running the corresponding small baselines under the same
  tokenizer, data budget, fine-tuning code, and seed protocol. Existing raw
  UniXcoder numbers are included as a placeholder reference.}
  \label{tab:comparison}
\end{table*}

\section{Results and Ablations}

Table~\ref{tab:comparison} shows the intended reporting format. The key claim
should be made only for models trained under the same scale and budget. Current
evidence suggests that raw CodeSearchNet pretraining is a strong baseline, while
Code-JEPA targets hard-negative separation. The projector-fixed Code-JEPA run
should replace older checkpoints trained before losses were moved from \(h\) to
\(z\).

\paragraph{Ablations.}
We report three compact ablations: (i) rank weight \(0\) versus \(0.5\), (ii)
directly training \(h\) versus training projected \(z\), and (iii) transformed
view stage or hard-negative source. The most important ablation is the projection
head, since downstream tasks consume \(h\).

\section{Limitations and Conclusion}

This work is not a reproduction of the original UniXcoder training recipe: it is
a small-model comparison built around Python CodeSearchNet. The full UniXcoder
recipe uses larger models, multiple languages, AST information, and additional
generation/denoising modes. The workshop version should therefore claim a
same-scale Code-JEPA comparison, not a full replacement for published foundation
models. Code-JEPA is most defensible as a method for adding hard-negative latent
prediction while preserving a reusable code embedding.

\end{document}
